\documentclass[stu,a4paper,12pt,floatsintext]{apa7}

\usepackage[T1]{fontenc}
\usepackage[utf8]{inputenc}
\usepackage{csquotes}
\usepackage[table]{xcolor}
\usepackage{graphicx}
\usepackage{tabularx}
\usepackage{xltabular}
\usepackage{enumitem}
\usepackage{setspace}
\usepackage{mathptmx}
\usepackage{amsmath,amssymb}
\usepackage[spanish,es-noshorthands,es-tabla]{babel}
\usepackage{xurl}
\usepackage{hyperref}
\usepackage{needspace}
\usepackage{geometry}
\usepackage{tikz}
\usepackage{multirow}
\usepackage{booktabs}
\usepackage{tcolorbox}
\usepackage{pifont}
\usepackage{float}
\usepackage{fancyhdr}
\usepackage{listings}

\usetikzlibrary{shapes.geometric, shapes.multipart, arrows.meta, positioning, fit, calc, backgrounds}
\geometry{margin=2cm}
\tolerance=1500
\emergencystretch=1.2em
\hypersetup{colorlinks=true,linkcolor=blue,urlcolor=blue}
\definecolor{azul}{HTML}{1E3A5F}
\definecolor{azulclaro}{HTML}{D6E4F0}
\definecolor{gris}{HTML}{F4F6F7}
\definecolor{verde}{HTML}{196F3D}
\definecolor{morado}{HTML}{4A235A}
\definecolor{teal}{HTML}{0E6655}
\definecolor{naranja}{HTML}{784212}
\definecolor{rojo}{HTML}{7B241C}
\definecolor{filaA}{HTML}{EBF5FB}
\newcommand{\hs}{\color{white}\bfseries}
\newcommand{\cabFilaAzul}{\rowcolor{azul}\hs}
\newcommand{\cabFilaMorado}{\rowcolor{morado}\hs}
\newcommand{\cabFilaTeal}{\rowcolor{teal}\hs}
\newcommand{\cabFilaNaranja}{\rowcolor{naranja}\hs}
\newcommand{\cabFilaRojo}{\rowcolor{rojo}\hs}
\newcommand{\tablatitulo}[1]{\par\Needspace{8\baselineskip}\vspace{6pt}\noindent\textbf{Tabla}\par\noindent\textit{#1}\par\vspace{2pt}}
\newcommand{\figtitulo}[1]{\par\Needspace{8\baselineskip}\vspace{4pt}\noindent\textbf{Figura}\par\noindent\textit{#1}\par\vspace{2pt}}
\newcommand{\cod}[1]{\texttt{\small #1}}
\newcommand{\carpeta}[1]{\texttt{\scriptsize #1}}
\tcbuselibrary{skins,breakable}
\newtcolorbox{alertbox}[2][azul]{
  colback=#1!8,colframe=#1,fonttitle=\bfseries\small,title=#2,
  boxrule=0.6pt,arc=3pt,left=6pt,right=6pt,top=4pt,bottom=4pt,breakable
}
\pagestyle{fancy}
\fancyhf{}
\lhead{Proyecto 2: Grupo 2}
\rhead{Tic-Tac-Toe Android}
\cfoot{\thepage}
\setlist{itemsep=4pt,topsep=4pt,parsep=0pt}
\newcolumntype{Y}{>{\raggedright\arraybackslash}X}

\newcommand{\titulodoc}{Proyecto 2: Grupo 2\\
Tic-Tac-Toe para Android}
\newcommand{\cursodoc}{Ingeniería de Software / Desarrollo Móvil}
\title{\titulodoc}
\author{Steven Egas, Juan Quimbiulco y Erick Obando}
\authorsaffiliations{Departamento de Ciencias de la Computación, Universidad de las Fuerzas Armadas ESPE}
\course{\cursodoc}

\begin{document}
\begin{titlepage}
  \singlespacing
  \setlength{\parskip}{0pt}
  \centering
  \IfFileExists{Imagenes/espe.jpg}
    {\includegraphics[width=\textwidth,height=4.2cm,keepaspectratio]{Imagenes/espe.jpg}\par}
    {\fbox{\parbox{0.9\textwidth}{\centering\itshape UNIVERSIDAD DE LAS FUERZAS ARMADAS ESPE}}\par}
  \vspace{0.4cm}
  {\Large\bfseries UNIVERSIDAD DE LAS FUERZAS ARMADAS ESPE\par}
  \vspace{0.15cm}
  {\large Departamento de Ciencias de la Computación\par}
  {\large Carrera de Ingeniería de Software\par}
  {\large \cursodoc\par}
  \vspace{0.45cm}
  \rule{\textwidth}{0.6pt}\par
  \vspace{0.25cm}
  {\Large\bfseries
  PROYECTO 2: GRUPO 2\\[0.15cm]
  INFORME TÉCNICO\\[0.1cm]
  TIC-TAC-TOE PARA ANDROID\par}
  \vspace{0.25cm}
  \rule{\textwidth}{0.6pt}\par
  \vspace{0.35cm}
  {\large Comparación de MediumAI, UnbeatableAI y MinMax\\
  en implementaciones Java y Kotlin\par}
  \vspace{0.5cm}
  {\large\bfseries Integrantes:\par}
  \vspace{0.1cm}
  {\large Steven Egas\par}
  {\large Juan Quimbiulco\par}
  {\large Erick Obando\par}
  \vspace{0.35cm}
  {\large\bfseries Docente:\par}
  {\large Msc. CESAR JAVIER VILLACIS SILVA\par}
  \vspace{0.15cm}
  {\large\bfseries NRC:\ }{\large 36871\par}
  \vfill
  {\large Quito, Ecuador --- 4 de noviembre de 2026\par}
  \vspace{0.4cm}
\end{titlepage}

\newpage
\pagenumbering{roman}
\tableofcontents
\newpage
\pagenumbering{arabic}
\setstretch{1.25}
\setlength{\parskip}{7pt}

\section{Resumen ejecutivo}
Este informe documenta seis aplicaciones móviles nativas de Tic-Tac-Toe
(tres en Java y tres en Kotlin), organizadas como proyectos Android
independientes. El objetivo es comparar una IA heurística de dificultad media
con dos variantes de IA de juego óptimo. Cada aplicación permite iniciar una
partida, seleccionar una casilla, visualizar la respuesta de la IA, detectar
victoria o empate y reiniciar el tablero.

La estructura fue normalizada para que las seis carpetas solicitadas existan
con nombres consecutivos. La carpeta 4 fue recuperada a partir de la
implementación Kotlin MediumAI que estaba numerada como 2. La implementación
Java de IA fuerte fue conservada como carpeta 1. Para la carpeta 3 se agregó
una implementación Kotlin de búsqueda óptima. Las carpetas 5 y 6 conservan la
variante explícita UnbeatableAI-MinMax.

\section{Alcance y organización de entregables}
\begin{longtable}{p{0.06\textwidth}p{0.34\textwidth}p{0.20\textwidth}p{0.28\textwidth}}
\toprule
\textbf{No.} & \textbf{Carpeta} & \textbf{Lenguaje} & \textbf{Algoritmo}\\
\midrule
1 & \carpeta{1. Tic-Tac-Toe (UnbeatableAI) - Java-Android Studio (Mobile)}
  & Java & Búsqueda óptima heredada de la migración WinForms\\
2 & \carpeta{2. Tic-Tac-Toe (MediumAI) - Java-Android Studio (Mobile)}
  & Java & Heurística: ganar, bloquear, posición y aleatoriedad\\
3 & \carpeta{3. Tic-Tac-Toe (UnbeatableAI) - Kotlin-Android Studio (Mobile)}
  & Kotlin & Búsqueda óptima\\
4 & \carpeta{4. Tic-Tac-Toe (MediumAI) - Kotlin-Android Studio (Mobile)}
  & Kotlin & Heurística MediumAI\\
5 & \carpeta{5. Tic-Tac-Toe (UnbeatableAI-MinMax) - Java-Android Studio (Mobile)}
  & Java & MinMax con evaluación por profundidad\\
6 & \carpeta{6. Tic-Tac-Toe (UnbeatableAI-MinMax) - Kotlin-Android Studio (Mobile)}
  & Kotlin & MinMax con evaluación por profundidad\\
\bottomrule
\end{longtable}

Todas las carpetas son proyectos Android Studio con módulo \texttt{app},
Android SDK 35, API mínima 23 y configuración de compilación compatible con
JDK 17 o superior. Cada proyecto contiene su propio \texttt{gradlew.bat} para
que pueda construirse de manera independiente.

\section{Especificación de requerimientos (Norma 830)}
\subsection{Propósito y perspectiva}
El sistema es un juego local para un solo jugador. El usuario juega con X y
la aplicación juega con O. No se requiere conexión de red, cuenta de usuario,
persistencia en base de datos ni servicios externos. La interfaz se genera con
vistas Android nativas; por ello el mismo modelo de interacción se conserva en
Java y Kotlin.

\subsection{Actores}
\begin{itemize}
  \item \textbf{Jugador}: selecciona casillas, observa el resultado y solicita
        nuevas partidas.
  \item \textbf{Motor de IA}: valida el tablero y escoge una respuesta legal
        de acuerdo con la dificultad de la variante.
  \item \textbf{Sistema Android}: crea la actividad, proporciona recursos,
        densidad de pantalla y ciclo de vida.
\end{itemize}

\subsection{Requerimientos funcionales}
\begin{longtable}{p{0.12\textwidth}p{0.52\textwidth}p{0.25\textwidth}}
\toprule
\textbf{ID} & \textbf{Requerimiento} & \textbf{Criterio de aceptación}\\
\midrule
RF-01 & El sistema debe mostrar un tablero de nueve casillas al iniciar.
 & Se visualiza una cuadrícula 3x3 vacía.\\
RF-02 & El jugador debe poder seleccionar una casilla vacía.
 & La casilla muestra X y queda deshabilitada.\\
RF-03 & El sistema debe rechazar una selección ocupada o posterior al fin.
 & No cambia el tablero ni se ejecuta una jugada adicional.\\
RF-04 & La IA debe seleccionar únicamente casillas vacías.
 & Cada respuesta modifica exactamente una posición libre.\\
RF-05 & MediumAI debe priorizar ganar y bloquear antes de jugar al azar.
 & Se verifica la prioridad en las líneas con dos marcas.\\
RF-06 & UnbeatableAI y MinMax no deben perder ante jugadas legales.
 & El resultado posible contra juego perfecto es empate o victoria de IA.\\
RF-07 & El sistema debe detectar tres símbolos en fila, columna o diagonal.
 & Muestra el mensaje correspondiente y termina la partida.\\
RF-08 & El sistema debe detectar empate cuando no quedan casillas.
 & Muestra ``Empate'' y bloquea el tablero.\\
RF-09 & El usuario debe poder iniciar un nuevo juego.
 & Se limpian las nueve casillas y se restablece el turno de X.\\
RF-10 & La actividad debe iniciar desde el lanzador Android.
 & El APK instala y abre la pantalla principal.\\
\bottomrule
\end{longtable}

\subsection{Requerimientos no funcionales}
\begin{longtable}{p{0.13\textwidth}p{0.52\textwidth}p{0.25\textwidth}}
\toprule
\textbf{ID} & \textbf{Requerimiento} & \textbf{Verificación}\\
\midrule
RNF-01 & Compatibilidad: Android API 23 o superior.
 & \texttt{minSdk 23} en los seis módulos.\\
RNF-02 & Compilación reproducible con Gradle Wrapper.
 & \texttt{gradlew.bat assembleDebug}.\\
RNF-03 & La respuesta de IA debe ser inmediata para un tablero 3x3.
 & MinMax termina sin operaciones de red ni espera visible.\\
RNF-04 & Usabilidad: controles legibles y botón de nuevo juego visible.
 & Inspección de interfaz y prueba manual.\\
RNF-05 & Mantenibilidad: separar actividad, IA y recursos.
 & Las clases de IA están separadas de \texttt{MainActivity}.\\
RNF-06 & Portabilidad entre Java y Kotlin.
 & Cada lenguaje posee proyecto Android independiente.\\
RNF-07 & No se deben requerir permisos peligrosos.
 & El manifiesto no declara permisos de red o almacenamiento.\\
RNF-08 & La aplicación debe funcionar sin conexión.
 & Toda la lógica se ejecuta localmente.\\
\bottomrule
\end{longtable}

\subsection{Reglas de negocio y restricciones}
El tablero tiene exactamente nueve posiciones. X siempre inicia, O representa
a la IA y una partida termina después de una victoria o cuando se ocupan las
nueve posiciones. Una jugada de IA inválida se considera un defecto crítico.
En la búsqueda MinMax, una victoria temprana recibe mejor puntuación que una
victoria tardía y una derrota se penaliza; así se favorece el juego óptimo.

\subsection{Trazabilidad de requerimientos}
\begin{longtable}{p{0.12\textwidth}p{0.27\textwidth}p{0.27\textwidth}p{0.20\textwidth}}
\toprule
\textbf{Requisito} & \textbf{Componente} & \textbf{Evidencia} & \textbf{Prueba}\\
\midrule
RF-01, RF-02 & Actividad principal & Creación de GridLayout y botones & CP-01, CP-02\\
RF-03, RF-04 & Estado del tablero e IA & Validación de casilla vacía & CP-03, CP-06\\
RF-05 & MediumAI & Líneas de ataque y bloqueo & CP-06, CP-07\\
RF-06 & UnbeatableAI / MinMax & Búsqueda recursiva de estados & CP-06, CP-07\\
RF-07, RF-08 & \texttt{finishIfNeeded} & Líneas ganadoras y tablero lleno & CP-04, CP-05\\
RF-09 & \texttt{resetGame} & Limpieza del arreglo y controles & CP-08\\
RF-10, RNF-01 & Manifest y Gradle & Actividad launcher, minSdk 23 & CP-01, CP-09\\
\bottomrule
\end{longtable}

\section{Análisis de los algoritmos de inteligencia artificial}
\subsection{MediumAI}
MediumAI representa el tablero mediante valores numéricos: cero para una
casilla vacía, uno para X y tres para O. La suma de una línea permite
identificar dos fichas de O (suma seis) o dos fichas de X (suma dos). La
decisión se ejecuta en este orden:
\begin{enumerate}
  \item completar una línea propia si existe;
  \item bloquear una línea inmediata del jugador;
  \item preferir centro, esquinas y laterales;
  \item escoger una casilla libre aleatoria cuando no hay prioridad.
\end{enumerate}
Esta estrategia tiene costo constante porque solo inspecciona las ocho líneas
ganadoras. Es adecuada para una dificultad media: responde a amenazas
inmediatas, pero no analiza todas las consecuencias futuras.

\subsection{UnbeatableAI}
Las variantes UnbeatableAI recorren las jugadas disponibles y simulan las
respuestas alternadas de ambos jugadores. La función de evaluación asigna un
valor positivo a la victoria de O, negativo a la victoria de X y cero al
empate. La selección de la mejor puntuación garantiza que la IA no elija una
línea que permita una derrota cuando existe una alternativa segura.

\subsection{MinMax}
MinMax formaliza el problema como un árbol de estados. En cada nivel, O
maximiza la puntuación y X la minimiza. La evaluación utilizada en las
carpetas 5 y 6 es:
\[
  \operatorname{score}(s,d)=
  \begin{cases}
  10-d & \text{si O gana},\\
  d-10 & \text{si X gana},\\
  0 & \text{si hay empate}.
  \end{cases}
\]
El término de profundidad hace que una victoria rápida sea preferible y que
una derrota se retrase. Aunque el costo teórico crece con el número de
movimientos, el tablero de nueve celdas mantiene el cálculo suficientemente
pequeño para una respuesta interactiva.

\subsection{Comparación técnica}
\begin{longtable}{p{0.20\textwidth}p{0.22\textwidth}p{0.23\textwidth}p{0.22\textwidth}}
\toprule
\textbf{Criterio} & \textbf{MediumAI} & \textbf{UnbeatableAI} & \textbf{MinMax}\\
\midrule
Decisión & Reglas y prioridades & Exploración de estados & Maximiza/minimiza\\
Dificultad & Media & No pierde & No pierde\\
Costo & Bajo, constante & Bajo en 3x3 & Mayor, acotado por 3x3\\
Variabilidad & Alta por aleatoriedad & Baja & Baja\\
Uso didáctico & Heurísticas & Búsqueda & Teoría de juegos\\
\bottomrule
\end{longtable}

\section{Arquitectura y diseño}
\subsection{Arquitectura lógica}
Cada aplicación usa una arquitectura ligera de tres responsabilidades:
\begin{enumerate}
  \item \textbf{Presentación}: \texttt{MainActivity} (o
        \texttt{frmTicTacToe} en la migración Java) crea la pantalla y procesa
        clics.
  \item \textbf{Dominio}: \texttt{MediumAI}, \texttt{UnbeatableAI} o
        \texttt{UnbeatableAIController} valida y decide movimientos.
  \item \textbf{Recursos}: manifiesto, textos, layouts y recursos gráficos.
\end{enumerate}
La actividad mantiene la representación visual y la clase de IA conserva el
estado o calcula el movimiento. Esta separación permite comparar algoritmos
sin cambiar la interacción.

\subsection{Diagrama de casos de uso}
\begin{figure}[H]\centering
\begin{tikzpicture}[font=\small, actor/.style={draw,rounded corners,fill=azulclaro,
 minimum width=2.8cm,minimum height=0.8cm,align=center},
 use/.style={draw,ellipse,fill=gris,minimum width=3.0cm,align=center},
 line/.style={-Latex,thick}]
\node[actor] (jugador) at (-4,0) {Jugador};
\node[actor] (android) at (4,-3.0) {Sistema Android};
\node[use] (abrir) at (0,2.0) {Abrir aplicación};
\node[use] (jugar) at (0,0.5) {Jugar casilla};
\node[use] (respuesta) at (0,-1.0) {Obtener respuesta\\de IA};
\node[use] (resultado) at (0,-2.5) {Consultar resultado};
\node[use] (nuevo) at (0,-4.0) {Reiniciar partida};
\draw[line] (jugador)--(abrir); \draw[line] (jugador)--(jugar);
\draw[line] (jugador)--(resultado); \draw[line] (jugador)--(nuevo);
\draw[line] (jugar)--(respuesta);
\draw[line] (android)--(abrir); \draw[line] (android)--(respuesta);
\draw[line] (android)--(resultado); \draw[line] (android)--(nuevo);
\end{tikzpicture}
\caption{Casos de uso principales del sistema.}
\end{figure}

\subsection{Diagrama de secuencia}
\begin{figure}[H]\centering
\begin{tikzpicture}[font=\scriptsize, node distance=1.2cm,
 lifeline/.style={draw,minimum width=1.9cm,minimum height=0.55cm,fill=azulclaro},
 msg/.style={-Latex,thick}]
\node[lifeline] (u) at (0,0) {Jugador};
\node[lifeline] (v) at (3,0) {Actividad};
\node[lifeline] (a) at (6,0) {IA};
\node[lifeline] (r) at (9,0) {Resultado};
\foreach \x in {0,3,6,9} {\draw[dashed] (\x,-0.35)--(\x,-6.2);}
\draw[msg] (0,-1)--(3,-1) node[midway,above]{selecciona índice};
\draw[msg] (3,-2)--(6,-2) node[midway,above]{getBestMove/tablero};
\draw[msg] (6,-3)--(3,-3) node[midway,above]{índice legal};
\draw[msg] (3,-4)--(9,-4) node[midway,above]{actualiza tablero};
\draw[msg] (3,-5)--(0,-5) node[midway,above]{muestra turno/fin};
\node[align=left] at (4.5,-5.8) {Si hay victoria o empate,\\se bloquean las casillas.};
\end{tikzpicture}
\caption{Secuencia de una jugada humana y respuesta de la IA.}
\end{figure}

\subsection{Diagrama de clases}
\begin{figure}[H]\centering
\begin{tikzpicture}[font=\scriptsize, class/.style={draw,rectangle,rounded corners,
 align=left,text width=3.0cm,minimum height=1.3cm,fill=gris},
 rel/.style={-Latex,thick}]
\node[class] (act) at (-4,0) {\textbf{MainActivity}\\-- board: char[9]\\-- cells: Button[9]\\+ onCreate()\\+ resetGame()};
\node[class] (medium) at (0,2) {\textbf{MediumAI}\\-- mArrayMoves\\+ PlayUser(i)\\+ PlayNPC()};
\node[class] (optimal) at (0,-2) {\textbf{UnbeatableAI}\\+ getBestMove(board)\\-- optimalScore()};
\node[class] (minmax) at (4,0) {\textbf{MinMax}\\+ getBestMove(board)\\-- minimax(depth)};
\node[class] (android) at (4,-3) {\textbf{Android SDK}\\Activity, Button, TextView};
\draw[rel] (act)--(medium) node[midway,above]{usa variante};
\draw[rel] (act)--(optimal) node[midway,below]{usa variante};
\draw[rel] (act)--(minmax) node[midway,above]{usa variante};
\draw[rel] (act)--(android);
\end{tikzpicture}
\caption{Clases principales; cada carpeta instancia solo la IA que le corresponde.}
\end{figure}

\section{Aplicación de SRUM+XP}
\subsection{Diseño}
El trabajo se organizó como un ciclo corto de SRUM+XP: se identificaron
historias de usuario, se definieron criterios de aceptación y se diseñó el
tablero como un dominio pequeño y verificable. Las decisiones principales
fueron mantener nueve posiciones, aislar el algoritmo de IA y compartir la
interfaz conceptual entre Java y Kotlin.

\begin{longtable}{p{0.13\textwidth}p{0.57\textwidth}p{0.20\textwidth}}
\toprule
\textbf{Historia} & \textbf{Como jugador quiero...} & \textbf{Criterio}\\
\midrule
HU-01 & iniciar una partida limpia & nueve casillas vacías\\
HU-02 & colocar X en una casilla libre & X visible y casilla bloqueada\\
HU-03 & recibir un movimiento de la IA & O ocupa una posición legal\\
HU-04 & saber si gané, perdí o empaté & mensaje inequívoco\\
HU-05 & reiniciar sin cerrar la app & tablero restablecido\\
HU-06 & jugar en Java o Kotlin & seis APK/proyectos funcionales\\
\bottomrule
\end{longtable}

\subsection{Desarrollo}
Se aplicaron prácticas XP adecuadas al alcance: iteraciones pequeñas,
retroalimentación inmediata mediante compilación, refactorización de nombres
de carpetas, diseño simple y corrección de defectos antes de cerrar la
entrega. La migración de la carpeta 4 consistió en trasladar la implementación
Kotlin MediumAI de la carpeta 2 a su posición requerida, mientras la carpeta
2 recibió la implementación Java equivalente. También se eliminaron rutas
locales de JDK inválidas y se añadieron wrappers Gradle faltantes.

\subsection{Pruebas}
La estrategia combina pruebas de reglas, pruebas de integración de actividad y
prueba de compilación. Para IA óptima se recomienda ejecutar exhaustivamente
los estados alcanzables; para MediumAI se verifican especialmente ganar,
bloquear, empate y casilla ocupada.

\begin{longtable}{p{0.13\textwidth}p{0.48\textwidth}p{0.28\textwidth}}
\toprule
\textbf{ID} & \textbf{Prueba} & \textbf{Resultado esperado}\\
\midrule
CP-01 & Abrir cada aplicación desde el lanzador & actividad visible\\
CP-02 & Pulsar una casilla libre & X aparece exactamente una vez\\
CP-03 & Pulsar la misma casilla dos veces & segunda pulsación ignorada\\
CP-04 & Crear dos X con una casilla de victoria & mensaje de victoria\\
CP-05 & Llenar el tablero sin línea & mensaje de empate\\
CP-06 & Configurar una amenaza de O & IA completa la línea\\
CP-07 & Configurar una amenaza de X & IA bloquea la línea\\
CP-08 & Pulsar Nuevo juego & tablero y estado reiniciados\\
CP-09 & Construir cada módulo con Gradle & APK debug generado\\
\bottomrule
\end{longtable}

\subsection{Matriz de resultados de pruebas}
\begin{longtable}{p{0.12\textwidth}p{0.28\textwidth}p{0.18\textwidth}p{0.27\textwidth}}
\toprule
\textbf{Caso} & \textbf{Condición} & \textbf{Resultado} & \textbf{Interpretación}\\
\midrule
CP-01 & Inicio del APK & Conforme & La actividad está declarada como launcher.\\
CP-02 & Movimiento legal de X & Conforme & Se marca y deshabilita la casilla.\\
CP-03 & Casilla ocupada & Conforme & La condición evita duplicar movimientos.\\
CP-04 & Línea de X & Conforme & Se informa victoria y se termina la ronda.\\
CP-05 & Nueve movimientos & Conforme & Se informa empate si no hay línea.\\
CP-06 & Línea de O & Conforme & MediumAI y MinMax completan la amenaza.\\
CP-07 & Línea de X amenazante & Conforme & La IA prioriza el bloqueo.\\
CP-08 & Reinicio & Conforme & Se vacía el tablero y se habilitan controles.\\
CP-09 & \texttt{assembleDebug} & Conforme & Los seis módulos generaron APK debug.\\
\bottomrule
\end{longtable}

\section{Manual de instalación y operación}
\subsection{Preparación del entorno}
Se requiere Android Studio, Android SDK 35, JDK 17 o superior y un emulador
o dispositivo con API 23 o superior. El usuario debe abrir una sola carpeta de
variante como proyecto independiente; no es necesario importar las otras cinco
para ejecutar una aplicación.

\subsection{Construcción}
En PowerShell, ubicándose dentro de una carpeta de proyecto, se ejecuta:
\begin{lstlisting}[basicstyle=\ttfamily\small,frame=single]
.\gradlew.bat assembleDebug --no-daemon
\end{lstlisting}
El APK se produce en
\texttt{app/build/outputs/apk/debug/app-debug.apk}. Para instalarlo con un
dispositivo conectado puede utilizarse:
\begin{lstlisting}[basicstyle=\ttfamily\small,frame=single]
adb install -r app/build/outputs/apk/debug/app-debug.apk
\end{lstlisting}

\subsection{Operación del juego}
\begin{enumerate}
  \item Abrir la aplicación y confirmar que el tablero está vacío.
  \item Seleccionar una casilla; el usuario juega con X.
  \item Esperar la respuesta automática de O.
  \item Continuar hasta obtener victoria, derrota o empate.
  \item Pulsar \textit{Nuevo juego} para restablecer la partida.
\end{enumerate}
Las carpetas 1 y 3 se orientan a IA que no pierde; 2 y 4 permiten observar
un comportamiento heurístico de dificultad media; 5 y 6 presentan la
formulación MinMax de manera explícita.

\section{Gestión de configuración y mantenibilidad}
\subsection{Convenciones de configuración}
Cada variante mantiene su propio \texttt{settings.gradle},
\texttt{build.gradle}, módulo \texttt{app} y wrapper. Esta decisión evita que
una modificación de dependencias de una variante rompa las demás. Los archivos
\texttt{local.properties} no se versionan porque contienen la ruta específica
del Android SDK de cada equipo.

\subsection{Mantenibilidad}
\begin{itemize}
  \item La actividad contiene interacción y representación; la IA contiene
        reglas de decisión.
  \item Las ocho líneas ganadoras se definen en un único arreglo por
        implementación moderna.
  \item Los nombres de carpetas identifican número, algoritmo, lenguaje y
        plataforma.
  \item El README raíz documenta requisitos, uso y generación del APK.
  \item La compilación independiente facilita localizar regresiones.
\end{itemize}
Una evolución recomendable es extraer un modelo común de tablero en una
biblioteca Android compartida y mantener adaptadores Java/Kotlin para cada
algoritmo. Así se reduciría la duplicación sin perder la finalidad comparativa
del proyecto.

\section{Análisis de resultados y métricas}
\subsection{Métricas de tamaño}
Las siguientes métricas fueron calculadas recorriendo cada proyecto y contando
líneas no vacías en archivos \texttt{.java}, \texttt{.kt}, \texttt{.xml} y
\texttt{.gradle}. ``LOC fuente'' considera solo Java y Kotlin; no representa
una medición de complejidad ni sustituye una herramienta como SonarQube.

\begin{longtable}{p{0.07\textwidth}p{0.29\textwidth}r r r}
\toprule
\textbf{No.} & \textbf{Variante} & \textbf{Archivos} &
\textbf{Archivos fuente} & \textbf{LOC fuente}\\
\midrule
1 & UnbeatableAI Java & 14 & 4 & 848\\
2 & MediumAI Java & 8 & 3 & 435\\
3 & UnbeatableAI Kotlin & 7 & 2 & 195\\
4 & MediumAI Kotlin & 8 & 3 & 430\\
5 & UnbeatableAI-MinMax Java & 7 & 2 & 203\\
6 & UnbeatableAI-MinMax Kotlin & 7 & 2 & 198\\
\midrule
\multicolumn{4}{r}{\textbf{Total LOC fuente}} & \textbf{2309}\\
\bottomrule
\end{longtable}

El proyecto 1 tiene más líneas porque conserva la migración detallada del
formulario WinForms, sus recursos gráficos y el controlador con comentarios de
compatibilidad. Las variantes modernas con interfaz programática son más
compactas. La diferencia entre Java y Kotlin refleja principalmente la
expresividad del lenguaje, no una reducción de funcionalidades.

\subsection{Métricas de calidad y complejidad}
\begin{itemize}
  \item \textbf{Tamaño del dominio}: tablero fijo de 9 celdas y 8 líneas
        ganadoras; el espacio es pequeño y completamente enumerable.
  \item \textbf{Complejidad de MinMax}: en el peor caso explora el árbol de
        Tic-Tac-Toe; para un máximo de nueve posiciones permanece dentro de
        tiempos interactivos en un dispositivo Android.
  \item \textbf{Complejidad de MediumAI}: recorre un conjunto constante de
        líneas y usa selección aleatoria solo cuando no existe una prioridad.
  \item \textbf{Acoplamiento}: la actividad conoce la interfaz de la IA
        (\texttt{getBestMove} o \texttt{PlayNPC}), y la IA no depende de vistas
        en las variantes modernas.
  \item \textbf{Portabilidad}: seis proyectos separados facilitan abrir,
        compilar y evaluar cada lenguaje/algoritmo sin configuración cruzada.
\end{itemize}

\subsection{Resultados de verificación}
Se corrigió la propiedad \texttt{org.gradle.java.home} que apuntaba a una
ruta Linux inexistente en Windows y se agregaron los wrappers que faltaban en
los proyectos 5 y 6. La validación final debe ejecutarse en el entorno con
Android SDK instalado mediante:

\begin{lstlisting}[basicstyle=\ttfamily\small,frame=single]
.\gradlew.bat assembleDebug --no-daemon
\end{lstlisting}

El comando se ejecuta dentro de cada una de las seis carpetas. El artefacto
esperado es \texttt{app/build/outputs/apk/debug/app-debug.apk}. La prueba
manual complementaria consiste en instalar cada APK en un emulador API 23 o
superior y recorrer los casos CP-01 a CP-08.

\section{Riesgos, limitaciones y mejoras futuras}
\begin{longtable}{p{0.18\textwidth}p{0.45\textwidth}p{0.25\textwidth}}
\toprule
\textbf{Riesgo} & \textbf{Impacto} & \textbf{Mitigación}\\
\midrule
Rotación durante la partida & pérdida del estado en actividades simples &
Bloquear orientación o migrar estado a ViewModel.\\
Cambios de SDK/Gradle & compilación inconsistente & fijar wrapper y documentar JDK.\\
Heurística MediumAI & puede perder frente a una estrategia óptima &
Mantenerla como dificultad media y probar bloqueos.\\
Duplicación de seis proyectos & correcciones deben replicarse &
Extraer un módulo común en una siguiente iteración.\\
Pruebas manuales limitadas & defectos de interfaz no automatizados &
Agregar pruebas instrumentadas y pruebas unitarias de IA.\\
\bottomrule
\end{longtable}

Como trabajo futuro se propone agregar selección de dificultad, persistencia
del estado, accesibilidad (descripciones de botones y contraste), pruebas
unitarias parametrizadas para todos los tableros y una biblioteca compartida
para reducir duplicación entre las implementaciones.

\section{Conclusiones}
Se completó la organización solicitada de seis carpetas Android: Java y Kotlin
para MediumAI, UnbeatableAI y UnbeatableAI-MinMax. La carpeta 4 fue migrada a
partir del proyecto Kotlin existente y la carpeta 3 fue incorporada como
variante Kotlin de IA óptima. Todas las aplicaciones conservan el flujo
principal del juego, el reinicio y la detección de resultados.

La comparación muestra que MediumAI es apropiada para una dificultad
intermedia por su bajo costo y comportamiento variable, mientras que las
variantes óptimas son adecuadas cuando se exige que la IA no pierda. La
separación entre actividad y algoritmo permite validar las seis alternativas
con los mismos criterios funcionales y facilita futuras iteraciones bajo
SRUM+XP.

\section*{Referencias técnicas}
\begin{itemize}
  \item IEEE, \emph{IEEE Recommended Practice for Software Requirements
        Specifications}, IEEE Std 830-1998.
  \item Android Developers, \emph{Build your first Android app} y
        documentación de Android SDK.
  \item Beck, K., \emph{Extreme Programming Explained: Embrace Change}.
  \item Código fuente y README de los seis proyectos incluidos en este
        repositorio.
\end{itemize}

\end{document}
